% ============================================
% LUC-16 — A Mãe no Telefone (edição revista R1075)
% ============================================
\fragtitle{luc}{L}{LUC}{16}{A Mãe no Telefone}
\markboth{LUC-16 --- A Mãe no Telefone}{}

\noindent \textit{Belo Horizonte --- 11 de janeiro de 2026, 22h17}

\noindent O telefone tocou três vezes antes de a mãe atender. Lúcia imaginou a casa inteira naquele intervalo: a sala com o sofá de flores, a mesa de crochê, o cheiro de lavanda e madeira velha, o retrato do pai na parede. A casa existia naquele cheiro. Fora dela, o pai virava uma foto; dentro, continuava sentado na poltrona, lendo o jornal com a ponta da língua entre os dentes.

\noindent ``Alô?''

\noindent ``Mãe. Sou eu.''

\noindent ``Lúcia.'' A voz da mãe mudou de tom. Não era surpresa. Era conferência. ``Aconteceu alguma coisa?''

\noindent ``Não. Só ligando.''

\noindent ``Você não liga só.''

\noindent Lúcia fechou os olhos. A mãe tinha razão. Ela ligava em datas marcadas, em aniversários, quando havia uma notícia prática: exame, consulta, uma reforma no telhado. Nunca só.

\noindent ``Estou bem'', disse. ``Estou trabalhando muito.''

\noindent ``Nesse arquivo?''

\noindent ``É.''

\noindent ``Aquele do hospital?''

\noindent Lúcia ouviu o barulho da televisão ao fundo. Novela. A mãe sempre deixava a TV ligada, mesmo sem assistir. Dizia que a casa ficava menos vazia.

\noindent ``É um caso antigo'', Lúcia disse. ``Estou organizando documentos.''

\noindent ``Você está com a voz rouca.''

\noindent ``É o ar-condicionado.''

\noindent ``Desde quando ar-condicionado deixa a voz rouca?'' A mãe fez uma pausa. ``Você está dormindo?''

\noindent Lúcia pensou em mentir. Não conseguiu.

\noindent ``Pouco.''

\noindent ``Eu sabia.''

\noindent O silêncio entre as duas era antigo, com regras próprias. A mãe esperava. Lúcia calculava. Havia coisas que ela nunca contava: o casamento que não aconteceu, o aborto aos trinta e dois anos, o namoro com Ronaldo, o medo de envelhecer sozinha. A mãe sabia de algumas. Não sabia de outras. As que ela não sabia eram as que mais pesavam.

\noindent ``A vizinha do 204 morreu'', a mãe disse.

\noindent ``Dona Cida?''

\noindent ``Ela mesma. Encontrada na cozinha. O cachorro latiu a noite inteira. Ninguém foi ver.''

\noindent ``Que triste.''

\noindent ``O cachorro está aqui no quintal. Ela deixou. Eu não queria, mas não tinha para quem dar.''

\noindent Lúcia ouviu um latido. Distante, abafado, como se viesse de dentro da parede.

\noindent ``Ele tem um olho amarelo'', a mãe disse. ``O direito. O esquerdo é marrom. Nunca vi cachorro assim.''

\noindent Lúcia abriu os olhos.

\noindent ``Mãe.''

\noindent ``Que é?''

\noindent ``Não deixa esse cachorro entrar em casa.''

\noindent A mãe riu. ``Você está com medo de cachorro agora?''

\noindent ``Medo, não. É...'' Lúcia parou. Não havia palavra. ``Só não deixa.''

\noindent Houve um estalo na linha. Um ruído fino, como folhas sendo viradas depressa. Depois, a voz da mãe voltou, mas havia algo errado com o volume. Estava mais baixa, mais lenta, como se falasse de dentro de um cômodo fechado.

\noindent ``Lúcia.''

\noindent ``Mãe?''

\noindent ``Você está com a voz do seu pai.''

\noindent Lúcia ficou parada. O pai morrera havia doze anos. Câncer de pulmão. Ele falava baixo, com pausas longas, como se cada frase custasse um pedaço de ar. Lúcia ouviu a própria respiração e percebeu que estava fazendo o mesmo.

\noindent ``Mãe, eu preciso desligar.''

\noindent ``Você não vai vir no domingo?''

\noindent ``Eu vou tentar.''

\noindent ``Não promete se não for cumprir.''

\noindent ``Eu prometo.''

\noindent Ela desligou antes que a mãe respondesse. Ficou com o telefone na mão, olhando para a tela apagada. Havia uma mancha de suor na capa. O cheiro de lavanda e madeira velha continuava no ar do apartamento, embora a mãe estivesse a duzentos quilômetros de distância. Lúcia levantou, abriu todas as janelas, deixou o vento entrar. O cheiro não saiu.

\noindent No celular, uma notificação: uma foto enviada pela mãe. O cachorro no quintal, sentado na grama, olhando para a câmera. O olho direito era amarelo. No canto da foto, atrás do cachorro, havia uma sombra com a forma de uma pessoa. Ninguém estava no quintal.

\noindent Lúcia olhou para a foto por muito tempo. Depois salvou, sem saber por quê, numa pasta chamada ``caso 1.260''. Era a primeira vez que ela usava o nome do paciente para uma imagem que não era documento.

\noindent\textit{\textbf{↪ Rotas:} \rota{LUC-15} \rotasep \rota{LUC-17} \rotasep \rota{MEM-07}}
